% apa7 loads hyperref internally; pass options before the class to avoid clash
\PassOptionsToPackage{hidelinks}{hyperref}
\documentclass[stu,12pt,floatsintext]{apa7}

% ── Encoding & font ────────────────────────────────────────────────────────────
\usepackage[T1]{fontenc}
\usepackage[utf8]{inputenc}
\usepackage{newtxtext,newtxmath}       % Times New Roman text and math

% ── Language & bibliography ────────────────────────────────────────────────────
\usepackage[american]{babel}
\usepackage{csquotes}
\usepackage[style=apa,sortcites=true,sorting=nyt,backend=biber]{biblatex}
\DeclareLanguageMapping{american}{american-apa}
\addbibresource{references.bib}

% ── Tables ─────────────────────────────────────────────────────────────────────
\usepackage{booktabs}
\usepackage{longtable}
\usepackage{array}

% ── Graphics & color ───────────────────────────────────────────────────────────
\usepackage{graphicx}
\usepackage{xcolor}

% ── Math ───────────────────────────────────────────────────────────────────────
\usepackage{amsmath}

% ── Figure placeholder box ─────────────────────────────────────────────────────
\newcommand{\figplaceholder}{%
  \begin{center}%
    \colorbox{gray!15}{%
      \parbox[c][3.6cm][c]{0.92\linewidth}{%
        \centering\small\itshape
        [Figure generated from analysis notebook --- replace with actual image file]%
      }%
    }%
  \end{center}%
}

% ── Document metadata ──────────────────────────────────────────────────────────
\title{Statistical Modeling and Air Quality Prediction Using Beijing Multi-Site Data}
\author{Ni Zhao, Pravin Suranthiran, and Sudhakaran Srinivasan}
\affiliation{Shiley-Marcos School of Engineering \\ University of San Diego}
\course{AAI-500: Applied Probability and Statistics for Artificial Intelligence}
\professor{Professor Leonid Shpaner}
\duedate{June 21, 2026}

\abstract{%
Air pollution is a major environmental and public health concern, especially in large urban
areas where pollutant concentrations vary by location, season, time of day, and
meteorological conditions. This project analyzes the Beijing Multi-Site Air Quality dataset
to identify urban air pollution patterns and evaluate whether air pollutant and
weather-related variables can be used to predict PM2.5 concentrations. The dataset
contains hourly observations from 12 monitoring stations in Beijing from March 2013
through February 2017. PM2.5 was selected as the primary target variable because it is a
widely used air-quality indicator and demonstrates meaningful variation across seasons and
locations. The project follows an end-to-end statistical analysis workflow that includes
data cleaning and preparation, exploratory data analysis, statistical hypothesis testing,
model selection, model analysis, and recommendations. Exploratory analysis shows that
PM2.5 is right-skewed, varies by station, is highest in winter and lowest in summer, and
has strong positive correlations with PM10, CO, NO2, and SO2. Statistical testing further
indicates that PM2.5 differs significantly across seasons and has a significant negative
relationship with wind speed. The model selection process compared linear regression,
random forest, logistic regression, naive Bayes, and XGBoost approaches across regression
and classification tasks. XGBoost was selected as the strongest overall model family, with
the best all-station forecast regression model producing $R^{2} = 0.515$,
RMSE $= 57.558$, and MAE $= 39.889$.%
}

\keywords{PM2.5, air quality, Beijing, exploratory data analysis, statistical modeling,
XGBoost}

% ══════════════════════════════════════════════════════════════════════════════
\begin{document}
\maketitle

\tableofcontents
\newpage

% ── Introduction ───────────────────────────────────────────────────────────────
\section{Introduction}

Air quality is an important environmental and public health issue. Fine particulate
matter, commonly referred to as PM2.5, is particularly important because small airborne
particles can enter deep into the respiratory system and affect health. The
U.S.\ Environmental Protection Agency \parencite{EPA2026} explains that small particles
can get deep into the lungs and, in some cases, enter the bloodstream. The World Health
Organization \parencite{WHO_nd} similarly describes air pollution as a major health
concern because exposure affects the lungs, heart, brain, and other organs.

The purpose of this project is to perform an end-to-end statistical analysis of the
Beijing Multi-Site Air Quality dataset. This dataset includes hourly measurements of six
major air pollutants and six meteorological variables from multiple monitoring sites in
Beijing \parencite{Chen2017}. The project objective is to understand patterns in PM2.5
concentrations and develop a statistical or machine learning model to predict PM2.5 based
on weather-related and pollutant variables.

This report is organized according to the required final project workflow. The data
cleaning/preparation section describes how the raw station files were combined and treated
for missing values. The exploratory data analysis section summarizes station-level,
seasonal, time-series, and correlation patterns. The statistical analysis section
evaluates whether observed seasonal differences and wind-speed relationships are
statistically significant. The model selection and model analysis sections compare
regression and classification approaches and explain why XGBoost was selected as the
strongest model family. The conclusion and recommendations section connects the technical
findings to environmental monitoring and decision-making.

% ── Data Cleaning/Preparation ──────────────────────────────────────────────────
\section{Data Cleaning/Preparation}

The Beijing Multi-Site Air Quality dataset was selected from the UCI Machine Learning
Repository. The UCI documentation describes the dataset as a multivariate time-series
dataset with 420,768 hourly observations, 12 monitoring stations, six major air pollutant
variables, and six relevant meteorological variables. The data period spans March~1, 2013
through February~28, 2017, and missing data are denoted as NA \parencite{Chen2017}.

The raw data were stored as 12 station-level CSV files. The first preparation step was to
load the CSV files, extract the station name from each file, and concatenate the records
into a single dataframe. This created a combined dataset with 420,768 observations and
18 original variables. A datetime feature was then created from the year, month, day, and
hour fields. After adding the datetime column, the working dataframe contained 19
variables.

\begin{table}[htbp]
  \begin{threeparttable}
  \caption{Selected Dataset Variables}
  \label{tab:variables}
  \begin{tabular}{ll}
    \toprule
    \textbf{Variable group}      & \textbf{Variables}                     \\
    \midrule
    Pollutants                   & PM2.5, PM10, SO2, NO2, CO, O3          \\
    Meteorological variables     & TEMP, PRES, DEWP, RAIN, wd, WSPM       \\
    Time variables               & year, month, day, hour, datetime        \\
    Location variable            & station                                 \\
    \bottomrule
  \end{tabular}
  \tablenote{The dataset contains hourly air pollution and meteorological measurements
    across 12 Beijing monitoring stations.}
  \end{threeparttable}
\end{table}

Missing values were reviewed across all columns. Pollutant variables had the largest
missing rates, with CO having the highest missing percentage. Weather variables had
smaller missing percentages. Table~\ref{tab:missing} summarizes the primary missing value
counts from the data cleaning notebook.

\begin{table}[htbp]
  \begin{threeparttable}
  \caption{Missing Value Summary Before Imputation}
  \label{tab:missing}
  \begin{tabular}{lcc}
    \toprule
    \textbf{Variable} & \textbf{Missing count} & \textbf{Missing \%} \\
    \midrule
    CO    & 20{,}701 & 4.92 \\
    O3    & 13{,}277 & 3.16 \\
    NO2   & 12{,}116 & 2.88 \\
    SO2   &  9{,}021 & 2.14 \\
    PM2.5 &  8{,}739 & 2.08 \\
    PM10  &  6{,}449 & 1.53 \\
    wd    &  1{,}822 & 0.43 \\
    WSPM  &    318   & 0.08 \\
    \bottomrule
  \end{tabular}
  \tablenote{Values were generated from the data-cleaning-preparation notebook.
    Only the most relevant missing variables are shown.}
  \end{threeparttable}
\end{table}

Numerical missing values were imputed using station-level medians. This choice was
appropriate because pollutant variables are right-skewed and contain extreme values, so
the median is less sensitive to unusually high pollution episodes than the mean. The
categorical wind direction variable was imputed using the station-level mode. After
imputation, the cleaned dataset had no remaining missing values and was saved as
\texttt{beijing\_air\_quality\_clean.csv} for additional analysis and modeling.
Imputation here runs the risk of data leakage, but it was done for the sake of dataset
consistency during model comparison.

Additional feature engineering was also performed. A season variable was created from the
month variable using the categories winter, spring, summer, and fall. This feature
supports seasonal comparison and hypothesis testing. For model preparation, hour and month
were also represented using sine and cosine transformations because time variables repeat
in cycles; for example, hour~23 and hour~0 are adjacent even though their raw numeric
values appear far apart.

% ── Exploratory Data Analysis ──────────────────────────────────────────────────
\section{Exploratory Data Analysis}

Exploratory data analysis was conducted to understand PM2.5 patterns before model
development. The EDA focused on PM2.5 distribution, station-level differences, seasonal
patterns, time-series trends, wind-speed relationships, and correlations among pollutant
and meteorological variables. This section connects directly to model selection by
identifying variables that may help predict PM2.5.

\subsection{PM2.5 Distribution}

The PM2.5 distribution is strongly right-skewed. Most observations are concentrated at
lower PM2.5 concentrations, while a smaller number of observations show very high
pollution values. The notebook summary statistics show a mean PM2.5 concentration of
approximately 79.79 and a median of 55.00. Since the mean is higher than the median, the
distribution demonstrates positive skewness. This matters for modeling because extreme
pollution events can influence error metrics and model fit.

\begin{figure}[htbp]
  \caption{Distribution of PM2.5 Concentration}
  \includegraphics[width=\linewidth]{figures/fig1.png}
  \label{fig:pm25dist}
  \figurenote{Most observations occur at lower PM2.5 levels, with a long right tail
    indicating occasional extreme pollution events.}
\end{figure}

\subsection{Station-Level PM2.5 Patterns}

Average PM2.5 levels varied across the 12 monitoring stations. The highest mean PM2.5
values in the current notebook output were observed at Dongsi, Wanshouxigong,
Nongzhanguan, and Gucheng. The lowest averages were observed at Dingling, Huairou, and
Changping. These differences suggest that PM2.5 concentrations are not evenly distributed
across Beijing monitoring locations. Station-level variation may reflect geographic
conditions, traffic density, local emission sources, and weather patterns.

\begin{table}[htbp]
  \begin{threeparttable}
  \caption{PM2.5 Summary Statistics by Station}
  \label{tab:station}
  \begin{tabular}{lccc}
    \toprule
    \textbf{Station}  & \textbf{Mean} & \textbf{Median} & \textbf{Standard deviation} \\
    \midrule
    Dongsi            & 86.19 & 61.00 & 86.58 \\
    Wanshouxigong     & 85.02 & 60.00 & 85.98 \\
    Nongzhanguan      & 84.84 & 59.00 & 86.23 \\
    Gucheng           & 83.85 & 60.00 & 82.80 \\
    Dingling          & 65.99 & 41.00 & 72.27 \\
    \bottomrule
  \end{tabular}
  \tablenote{Table shows selected high and low station-level averages from the data
    preparation notebook.}
  \end{threeparttable}
\end{table}

\begin{figure}[htbp]
  \caption{Average PM2.5 Concentration by Monitoring Station}
  \includegraphics[width=\linewidth]{figures/fig2.png}
  \label{fig:station}
  \figurenote{The station-level bar chart shows location-based variation in PM2.5
    concentrations.}
\end{figure}

\subsection{Seasonal PM2.5 Patterns}

Seasonal analysis showed that PM2.5 concentrations were highest in winter and lowest in
summer. The seasonal mean PM2.5 concentration was approximately 95.48 in winter, 82.33
in fall, 76.97 in spring, and 64.67 in summer. This pattern suggests that meteorological
conditions and seasonal activity patterns may influence PM2.5 levels. Winter pollution
may be related to lower atmospheric mixing, heating demand, and weather conditions that
trap pollutants near the ground.

\begin{table}[htbp]
  \begin{threeparttable}
  \caption{PM2.5 Summary Statistics by Season}
  \label{tab:season}
  \begin{tabular}{lcccc}
    \toprule
    \textbf{Season} & \textbf{Count}  & \textbf{Mean} & \textbf{Median}
                    & \textbf{Standard deviation} \\
    \midrule
    Winter & 101{,}907 & 95.48 & 56.00 & 107.66 \\
    Fall   & 102{,}746 & 82.33 & 56.00 &  81.44 \\
    Spring & 103{,}705 & 76.97 & 58.00 &  68.83 \\
    Summer & 103{,}671 & 64.67 & 51.00 &  52.65 \\
    \bottomrule
  \end{tabular}
  \tablenote{Season was derived from month: winter = December--February,
    spring = March--May, summer = June--August, and fall = September--November.}
  \end{threeparttable}
\end{table}

\begin{figure}[htbp]
  \caption{Average PM2.5 by Season}
  \includegraphics[width=\linewidth]{figures/fig3.png}
  \label{fig:seasonmean}
  \figurenote{Winter had the highest average PM2.5 concentration, while summer had
    the lowest average PM2.5 concentration.}
\end{figure}

\begin{figure}[htbp]
  \caption{PM2.5 Distribution by Season}
  \includegraphics[width=\linewidth]{figures/fig4.png}
  \label{fig:seasonbox}
  \figurenote{The boxplot shows that PM2.5 varies by season and contains high outlier
    values in each season.}
\end{figure}

\subsection{Time-Series PM2.5 Patterns}

Monthly trends by year further support the presence of recurring seasonal patterns. To
support fair year-over-year comparison, the notebook used complete years only: 2014,
2015, and 2016. The resulting line plot shows recurring winter peaks and summer reductions
in PM2.5. The pattern appears more seasonal than year-specific, suggesting that seasonal
timing is an important factor for predicting PM2.5.

\begin{figure}[htbp]
  \caption{Monthly PM2.5 Trend by Year}
  \includegraphics[width=\linewidth]{figures/fig5.png}
  \label{fig:trend}
  \figurenote{The chart uses full years only and shows recurring winter peaks and summer
    decreases across multiple years.}
\end{figure}

\subsection{Correlation and Wind-Speed Relationships}

Correlation analysis was used to identify variables that may be important for modeling
PM2.5. PM2.5 had strong positive relationships with PM10 ($r = 0.88$), CO ($r = 0.79$),
and NO2 ($r = 0.67$). Wind speed had a negative relationship with PM2.5 ($r = -0.27$).
These findings suggest that pollutant variables are strongly related to each other, while
stronger winds are associated with lower PM2.5 concentrations.

\begin{figure}[htbp]
  \includegraphics[width=\linewidth]{figures/fig6.png}
  \caption{Correlation Heatmap of Air Quality and Meteorological Variables}
  \label{fig:heatmap}
  \figurenote{PM2.5 is strongly correlated with other pollutant variables and negatively
    correlated with wind speed.}
\end{figure}

The scatterplot of PM2.5 and wind speed shows that severe PM2.5 episodes occur primarily
during low wind conditions. This pattern supports the idea that wind can help disperse
pollutants and reduce the likelihood of extreme particulate matter concentration.

\begin{figure}[htbp]
  \includegraphics[width=\linewidth]{figures/fig7.png}
  \caption{PM2.5 and Wind Speed Relationship}
  \label{fig:windscatter}
  \figurenote{Extreme PM2.5 values appear more common when wind speed is low.}
\end{figure}

% ── Statistical Analysis ───────────────────────────────────────────────────────
\section{Statistical Analysis}

Inferential statistical tests were conducted to determine whether observed EDA patterns
were statistically meaningful. These tests help connect descriptive findings to
probability and statistics concepts covered in the course.

\subsection{Seasonal Differences in PM2.5}

The first hypothesis test examined whether PM2.5 concentrations differed significantly
across seasons. Because there are four seasonal groups, a one-way analysis of variance
(ANOVA) was used.

\textit{Null hypothesis}: Mean PM2.5 concentration is the same across winter, spring,
summer, and fall.

\textit{Alternative hypothesis}: At least one season has a different mean PM2.5
concentration.

The ANOVA test produced $F = 2617.12$ and $p < .001$. Therefore, the null hypothesis
was rejected. This result indicates that PM2.5 concentrations differ significantly across
seasons. This supports the EDA finding that winter has the highest average PM2.5 and
summer has the lowest average PM2.5.

\subsection{Association Between Wind Speed and PM2.5}

The second hypothesis test examined the relationship between wind speed and PM2.5
concentration. A Pearson correlation significance test was used because both variables
are numeric.

\textit{Null hypothesis}: Wind speed has no significant relationship with PM2.5
concentration.

\textit{Alternative hypothesis}: Wind speed is significantly associated with PM2.5
concentration.

The Pearson correlation test produced $r = -0.27$ and $p < .001$. Therefore, the null
hypothesis was rejected. The negative correlation indicates that higher wind speeds are
generally associated with lower PM2.5 concentrations. Although the relationship is not
extremely strong, it is statistically significant and directionally meaningful for air
pollution interpretation.

% ── Model Selection ────────────────────────────────────────────────────────────
\section{Model Selection}

The modeling objective was to predict PM2.5 concentration or PM2.5 pollution category
using pollutant, meteorological, and time-related variables. PM2.5 was selected as the
target variable because it is an important air-quality indicator and because the EDA
demonstrated strong variation by station, season, year, and weather conditions.

Two prediction settings were evaluated. The first setting was a nowcasting task, which
predicted PM2.5 using concurrent pollutant and weather readings, including SO2, NO2, CO,
O3, temperature, pressure, dew point, rainfall, and wind speed. The second setting was a
forecasting-style task, which predicted PM2.5 using weather and time features only,
without concurrent pollutant measurements. Comparing these settings is useful because
nowcasting represents the best available information scenario, while the weather-only
setting is closer to a future prediction problem.

For model comparison, PM10 was excluded because it was highly correlated with PM2.5 and
could dominate the model. Wind direction was also excluded from the model-selection
notebook. Rows with missing values were dropped in the modeling dataframe so that all
models were evaluated on the same complete-case dataset. After this step, the
model-selection dataframe contained 383,767 observations. Hour and month were encoded as
sine and cosine features to preserve cyclical time patterns. An 80/20 temporal split was
used, meaning the earlier 80\% of observations formed the training set and the final
20\% formed the test set. This avoided shuffling future observations into training and
reduced the risk of temporal data leakage.

The model-selection notebook evaluated both regression and classification formulations.
For regression, PM2.5 remained a continuous numeric target and models were evaluated
using root mean squared error (RMSE), mean absolute error (MAE), and $R$-squared. For
classification, PM2.5 was discretized into quintile classes and binary high/low classes
using the training distribution. Classification models were evaluated using accuracy and
weighted $F_1$ score. The models evaluated included ordinary least squares linear
regression, random forest regression/classification, logistic regression, Gaussian naive
Bayes, and XGBoost. XGBoost is a gradient-boosted tree ensemble that sequentially fits
trees to reduce prior model errors and can capture nonlinear relationships and
interactions \parencite{ChenGuestrin2016}.

\begin{table}[htbp]
  \begin{threeparttable}
  \caption{Regression Model Performance for All Stations}
  \label{tab:regression}
  \begin{tabular}{llccc}
    \toprule
    \textbf{Scenario} & \textbf{Model} & \textbf{RMSE} & \textbf{MAE}
                      & $\boldsymbol{R^{2}}$ \\
    \midrule
    Nowcast: pollutants + weather & XGBoost       & 38.945 & 23.395 & 0.778 \\
    Nowcast: pollutants + weather & Random Forest & 38.666 & 23.426 & 0.781 \\
    Nowcast: pollutants + weather & OLS           & 42.079 & 27.956 & 0.741 \\
    Forecast: weather only        & XGBoost       & 57.558 & 39.889 & 0.515 \\
    Forecast: weather only        & Random Forest & 63.416 & 43.562 & 0.411 \\
    Forecast: weather only        & OLS           & 65.754 & 46.713 & 0.367 \\
    \bottomrule
  \end{tabular}
  \tablenote{Lower RMSE and MAE indicate lower prediction error. Higher $R^{2}$
    indicates a larger proportion of PM2.5 variance explained.}
  \end{threeparttable}
\end{table}

For the all-station regression task, XGBoost performed best in both scenarios. In the
nowcast setting, XGBoost achieved $R^{2} = 0.778$, RMSE $= 38.945$, and MAE $= 23.395$.
In the weather-only forecast setting, performance declined to $R^{2} = 0.515$,
RMSE $= 57.558$, and MAE $= 39.889$. This decline is expected because the forecast
setting removes concurrent pollutant measurements that provide strong direct signals about
air quality.

\begin{figure}[htbp]
  \caption{Regression Results: $R$-Squared Comparison}
  \includegraphics[width=\linewidth]{figures/fig8.png}
  \label{fig:rsq}
  \figurenote{XGBoost produced the highest $R$-squared in both the nowcast and
    weather-only forecast scenarios.}
\end{figure}

\begin{table}[htbp]
  \begin{threeparttable}
  \caption{XGBoost Classification Performance for All Stations}
  \label{tab:class}
  \begin{tabular}{llcc}
    \toprule
    \textbf{Scenario} & \textbf{Classification task} & \textbf{Accuracy}
                      & \textbf{Weighted $F_1$} \\
    \midrule
    Nowcast: pollutants + weather & PM2.5 quintile & 0.616 & 0.613 \\
    Nowcast: pollutants + weather & PM2.5 binary   & 0.870 & 0.870 \\
    Forecast: weather only        & PM2.5 quintile & 0.456 & 0.434 \\
    Forecast: weather only        & PM2.5 binary   & 0.768 & 0.768 \\
    \bottomrule
  \end{tabular}
  \tablenote{Binary classification separates PM2.5 above and below the training median.
    Quintile classification separates PM2.5 into five ordered groups.}
  \end{threeparttable}
\end{table}

Classification results followed the same overall pattern as regression. Binary
classification was much easier than quintile classification because it only required the
model to separate higher and lower PM2.5 observations. XGBoost achieved 0.870 accuracy
and 0.870 weighted $F_1$ in the nowcast binary task and 0.768 accuracy and 0.768 weighted
$F_1$ in the weather-only binary task. Quintile classification was more difficult,
especially in the middle bins where PM2.5 values are less clearly separated.

The model-selection analysis also compared pooled all-station models against single-station
models for each monitoring site. XGBoost was strictly superior in all forecasting tasks and 
outperformed in a majority of nowcasting tasks as well. Notably, for certain highly linear stations, OLS provided
comparable and sometimes superior results, showing the weakness of tree models for highly linear
datasets and those with outlier data which OLS can effectively extrapolate. 
For some monitoring sites, Random Forest outperformed XGBoost, likely due to the tendency
of XGBoost to overfit to noise for simpler datasets.
In most cases, the pooled model performed similarly or better than the corresponding single-station model,
indicating that a general model across sites is viable. For instance, for the Aotizhongxin station, 
the pooled model achieved an $R^{2}$ of 0.818 in the nowcast setting, while the single-station model achieved 0.820; however,
for other stations, the pooled model often yielded higher $R^{2}$ values. 
One station-level result is provided in the appendix and comprehensive station-level analysis is included in the notebook.

% ── Model Analysis ─────────────────────────────────────────────────────────────
\section{Model Analysis}

XGBoost was selected as the final model family because it showed the strongest and most
consistent performance across regression, quintile classification, and binary
classification tasks. It also performed well in both the nowcast and forecast-style
settings. This consistency is important because a strong final model should not only
perform well under one narrow scenario; it should also remain competitive when features
are limited or when the target is reframed as a classification problem.

The model results show a clear distinction between nowcasting and forecasting. The nowcast
model performed best because concurrent pollutant measurements such as CO, NO2, SO2, and
O3 provide strong direct information about current PM2.5 conditions. The weather-only
model performed worse but still explained approximately half of the PM2.5 variation in the
all-station regression task. This suggests that meteorological and seasonal conditions are
meaningful predictors, but weather alone cannot fully identify emission events or
short-term pollution spikes.

\begin{figure}[htbp]
  \caption{XGBoost Actual Versus Predicted PM2.5 Across Four Scenarios}
  \includegraphics[width=\linewidth]{figures/fig9.png}
  \label{fig:actualvpred}
  \figurenote{Predictions closer to the diagonal line indicate better model fit.
    Nowcast scenarios show stronger alignment than weather-only forecast scenarios.}
\end{figure}

Feature-importance output further supports the EDA findings. In the nowcast scenario,
pollutant features, especially CO and NO2, contributed strongly to model performance. In
the weather-only forecast scenario, temporal features, dew point, pressure, and other
meteorological variables became more important. This suggests that the nowcast model
relies on direct pollutant signals, while the forecast-style model relies more heavily on
seasonal and weather patterns.

\begin{figure}[htbp]
  \caption{XGBoost Feature Importance}
  \includegraphics[width=\linewidth]{figures/fig10.png}
  \label{fig:importance}
  \figurenote{The nowcast model is driven more by pollutant variables, while the
    weather-only model relies more on time and meteorological variables.}
\end{figure}

The selected model is valid for the course project because it is evaluated using a
temporal train/test split, compared against simpler baseline models, and interpreted using
multiple metrics. However, there are important limitations. First, the nowcast model
should not be interpreted as a true future forecast because it uses concurrent pollutant
measurements. Second, the weather-only forecast model performs less strongly and may miss
sudden emission events. Third, PM10 and lagged PM2.5 were excluded to prevent them from
dominating the model comparison, but these variables could improve predictive performance
in a real operational system. Fourth, the model does not include external predictors such
as traffic volume, industrial activity, policy changes, holidays, or regional weather
systems.

Overall, XGBoost provides the best balance of predictive performance and flexibility for
this dataset. It can capture nonlinear relationships, tolerate correlated predictors
better than linear models, and provide feature-importance outputs that support
interpretation. The most appropriate final conclusion is that XGBoost is the strongest
model family for PM2.5 prediction in this project, with the nowcast version performing
best when concurrent pollutant readings are available and the weather-only version
offering a more realistic but less precise forecasting baseline.

% ── Conclusion and Recommendations ────────────────────────────────────────────
\section{Conclusion and Recommendations}

This project analyzed the Beijing Multi-Site Air Quality dataset to understand PM2.5
pollution patterns and develop a predictive modeling workflow. The data cleaning process
combined 12 station-level files, created a datetime feature, engineered seasonal
groupings, and treated missing values using station-level median and mode imputation. The
cleaned dataset supported exploratory analysis, statistical testing, and model
development.

The EDA showed that PM2.5 is right-skewed, varies across monitoring stations, is highest
in winter and lowest in summer, and has recurring seasonal patterns across years.
Correlation analysis showed strong positive relationships between PM2.5 and PM10, CO, and
NO2, while wind speed had a significant negative relationship with PM2.5. Statistical
testing confirmed that PM2.5 differs significantly across seasons and that wind speed is
significantly associated with PM2.5.

The model-selection results support XGBoost as the final model family. In the all-station
nowcast regression task, XGBoost achieved $R^{2} = 0.778$, RMSE $= 38.945$, and
MAE $= 23.395$. In the weather-only forecast-style task, XGBoost remained the best model
but declined to $R^{2} = 0.515$. This gap is meaningful: it shows that concurrent
pollutant readings substantially improve current PM2.5 estimation, while weather and time
variables provide useful but incomplete predictive information.

From a practical perspective, this type of analysis can support environmental monitoring,
public health planning, and early warning systems for poor air quality. A nowcast model
could support real-time air-quality estimation when multiple pollutant measurements are
available. A weather-only model could serve as a baseline for forecasting periods when
pollutant data are unavailable or delayed. Future work should add lagged PM2.5, traffic
activity, industrial activity, policy changes, holidays, and additional regional weather
measures. Future model evaluation should also focus more specifically on high-pollution
tail events, because correctly identifying severe pollution episodes may be more important
for public health decision-making than minimizing average prediction error.

% ── References ─────────────────────────────────────────────────────────────────
\nocite{*}
\printbibliography

% ══════════════════════════════════════════════════════════════════════════════
% Appendices
% ══════════════════════════════════════════════════════════════════════════════
\appendix

% ── Appendix A ─────────────────────────────────────────────────────────────────
\section{Technical Notebook Output Summary}

The appendix summarizes key outputs from the project notebooks. The full notebooks are
maintained in the team GitHub repository. The report body includes the main figures and
tables needed to support the analysis.

\begin{table}[htbp]
  \begin{threeparttable}
  \caption{Key Technical Outputs Included in the Current Draft}
  \label{tab:outputs}
  \begin{tabular}{p{3cm}p{9.5cm}}
    \toprule
    \textbf{Notebook area}   & \textbf{Current output included} \\
    \midrule
    Data loading             & 12 station CSV files loaded and combined \\
    Combined dataset         & 420,768 observations; 12 stations; March 2013--February 2017 \\
    Missing values           & Missing count and percentage by variable \\
    Feature engineering      & datetime, season, \texttt{hour\_sin}/\texttt{hour\_cos},
                               and \texttt{month\_sin}/\texttt{month\_cos} columns created \\
    EDA figures              & PM2.5 distribution, station chart, seasonal charts,
                               correlation heatmap, wind-speed scatterplot, monthly trend \\
    Statistical tests        & ANOVA for seasonal PM2.5 differences; Pearson correlation
                               for wind speed and PM2.5 \\
    Modeling                 & Regression, quintile classification, and binary
                               classification model comparisons \\
    \bottomrule
  \end{tabular}
  \tablenote{Outputs are summarized from the data-cleaning-preparation,
    \texttt{nina\_eda\_air\_quality}, and model-selection notebooks.}
  \end{threeparttable}
\end{table}

% ── Appendix B ─────────────────────────────────────────────────────────────────
\section{Complete Model Comparison Output}

\begin{table}[htbp]
  \begin{threeparttable}
  \caption{Regression Performance Across All Model Scenarios}
  \label{tab:regfull}
  \small
  \begin{tabular}{>{}p{3.0cm}>{}p{2.5cm}>{}p{2.5cm}ccc}
    \toprule
    \textbf{Feature set} & \textbf{Scope} & \textbf{Model}
      & \textbf{RMSE} & \textbf{MAE} & $\boldsymbol{R^{2}}$ \\
    \midrule
    Pollutants + weather & All stations & Random Forest & 38.666 & 23.426 & 0.781 \\
    Pollutants + weather & All stations & XGBoost       & 38.945 & 23.395 & 0.778 \\
    Pollutants + weather & All stations & OLS           & 42.079 & 27.956 & 0.741 \\
    Pollutants + weather & Aotizhongxin & XGBoost       & 35.390 & 22.274 & 0.820 \\
    Pollutants + weather & Aotizhongxin & OLS           & 38.857 & 26.141 & 0.783 \\
    Pollutants + weather & Aotizhongxin & Random Forest & 40.723 & 24.799 & 0.762 \\
    Weather only         & All stations & XGBoost       & 57.558 & 39.889 & 0.515 \\
    Weather only         & All stations & Random Forest & 63.416 & 43.562 & 0.411 \\
    Weather only         & All stations & OLS           & 65.754 & 46.713 & 0.367 \\
    Weather only         & Aotizhongxin & XGBoost       & 58.036 & 39.994 & 0.516 \\
    Weather only         & Aotizhongxin & Random Forest & 63.200 & 43.952 & 0.426 \\
    Weather only         & Aotizhongxin & OLS           & 65.320 & 46.220 & 0.387 \\
    \bottomrule
  \end{tabular}
  \tablenote{Model-selection notebook output; feature set and scope combinations used
    an 80/20 temporal train/test split.}
  \end{threeparttable}
\end{table}

\begin{table}[htbp]
  \begin{threeparttable}
  \caption{Classification Performance Across All-Stations Model Scenarios}
  \label{tab:classfull}
  \footnotesize
  \begin{tabular}{p{3.0cm}p{1.8cm}p{3.2cm}cc}
    \toprule
    \textbf{Feature set} & \textbf{Task} & \textbf{Model}
      & \textbf{Accuracy} & \textbf{Weighted $F_1$} \\
    \midrule
    Pollutants + weather & Quintile & XGBoost             & 0.616 & 0.613 \\
    Pollutants + weather & Quintile & Logistic Regression & 0.602 & 0.601 \\
    Pollutants + weather & Quintile & Random Forest       & 0.567 & 0.556 \\
    Pollutants + weather & Quintile & Naive Bayes         & 0.483 & 0.462 \\
    Pollutants + weather & Binary   & XGBoost             & 0.870 & 0.870 \\
    Pollutants + weather & Binary   & Logistic Regression & 0.857 & 0.856 \\
    Pollutants + weather & Binary   & Random Forest       & 0.845 & 0.845 \\
    Pollutants + weather & Binary   & Naive Bayes         & 0.770 & 0.764 \\
    Weather only         & Quintile & XGBoost             & 0.456 & 0.434 \\
    Weather only         & Quintile & Logistic Regression & 0.443 & 0.425 \\
    Weather only         & Quintile & Random Forest       & 0.395 & 0.372 \\
    Weather only         & Quintile & Naive Bayes         & 0.341 & 0.320 \\
    Weather only         & Binary   & XGBoost             & 0.768 & 0.768 \\
    Weather only         & Binary   & Logistic Regression & 0.729 & 0.728 \\
    Weather only         & Binary   & Random Forest       & 0.632 & 0.617 \\
    Weather only         & Binary   & Naive Bayes         & 0.615 & 0.613 \\
    \bottomrule
  \end{tabular}
  \tablenote{Classification results shown for pooled all-station experiments only.
    The binary task separates PM2.5 into below/above median categories based on
    training data.}
  \end{threeparttable}
\end{table}

% ── Appendix C ─────────────────────────────────────────────────────────────────
\section{Additional Model Figures}

\begin{figure}[htbp]
  \caption{Classification Results: Accuracy}
  \includegraphics[width=\linewidth]{figures/fig11.png}
  \label{fig:classacc}
  \figurenote{Quintile classification is harder than binary classification,
    particularly in the weather-only forecast setting.}
\end{figure}

\begin{figure}[htbp]
  \caption{Binary Classification Results: Accuracy}
  \includegraphics[width=\linewidth]{figures/fig12.png}
  \label{fig:binacc}
  \figurenote{Binary classification accuracy is substantially above the 50\% chance
    baseline for all model scenarios.}
\end{figure}

\begin{figure}[htbp]
  \caption{XGBoost Quintile Classification Confusion Matrix}
  \includegraphics[width=\linewidth]{figures/fig13.png}
  \label{fig:qcm}
  \figurenote{Extreme quintiles are generally easier to classify than middle
    quintiles.}
\end{figure}

\begin{figure}[htbp]
  \caption{XGBoost Binary Classification Confusion Matrix}
  \label{fig:bcm}
  \includegraphics[width=\linewidth]{figures/fig14.png}
  \figurenote{The binary confusion matrix shows stronger class separation than the
    five-quintile task.}
\end{figure}

% ── Appendix D ─────────────────────────────────────────────────────────────────
\section{Code and Repository Information}

GitHub repository:
\url{https://github.com/sudhakaran-srinivasan/air-quality-statistical-modeling}

Relevant notebooks in the repository include
\texttt{final\_modeling.ipynb},
\texttt{data-cleaning-preparation.ipynb},
\texttt{data-exploration.ipynb},
\texttt{nina\_eda\_air\_quality.ipynb}, 
and
\texttt{model-selection.ipynb}.

\end{document}